\section*{5. Prototype}
\addcontentsline{toc}{section}{5. Prototype}

\subsection*{Definition}
\addcontentsline{toc}{subsection}{Definition}
etc. etc.

\subsection*{Personal Reflection}
My primary role in the project was to document and write up about an example of an app for the project, creating user scenarios and user stories. Everyone 
in the group separated the work equally and sometimes there would be two people working on the same section to help each other out.
I also contributed to making the UML diagrams such as the class diagram. I've also teamed up with one of the group members to help create the activity diagram as well. I've learned how to use Visual Paradigm to create the diagrams. 
By using the lecturer's example of the UML diagrams, I was able to understand how to design the diagrams and how it is represented in Visual Paradigm. I found it to be a useful tool and it has a lot of features that helped me what the shapes and symbols mean. 
It was a bit confusing at the start since it has a lot of options, but once I got used to it, it was easy to use.
Creating the user scenarios and user stories were fun to do as it allowed me to be creative and think about how the app would be used in real life.

Our biggest challenge was keeping the group on track and keeping everyone consisten with the work.
We addressed this by chatting with each other on a regular basis on Discord and making sure that we get
the work done before the deadline. It was difficult to keep on track as everyone has different schedules and there would be exams and other assignments to do, 
but we managed to get the work done on time. 

Choosing Latex as our documentation tool was a learning experience.For myself, Latex was a new tool that I rarely used before, especially using it in VSCode while connected to GitHub. 
The start-up was difficult as I was unfamiliar with the syntax and spent a lot of time trying to figure out things. But one of my group member had experience with Latex and helped me out with the syntax and how to use it.
Once I got passed the learning curve, I found it to be very useful and professional for documentation. It allowed us to create a well-structured and formatted document that looks great. I also learned how to use GitHub to 
collaborate with my group members and keep track of the changes we made to the document.

In the future project I would like to improve my time management skills and communication with my group members. I think that would help us to keep on track and make sure 
that everyone is on the same page with the work. I also want to learn more about Latex and how to use it more efficiently for documentation. 
Overall, I enjoyed working on this project and learned a lot from it.


\subsection*{Tools Used}
\addcontentsline{toc}{subsection}{Tools Used}
 description, ease of use, brief.
\subsection*{Media (URLs, Screenshots)}
\addcontentsline{toc}{subsection}{Media (URLs, Screenshots)}
 sample of relevant urls, screenshots.

 Studo Website :
 https://studo.com/en

 Studo Reviews :
 https://chrome-stats.com/d/com.moshbit.studo
 